\documentclass[11pt]{article}
\usepackage{amsmath, amssymb, amsthm}
\usepackage{geometry}
\usepackage{hyperref}
\usepackage{array}
\usepackage{booktabs}
\usepackage{longtable}
\geometry{margin=1in}

\newtheorem{theorem}{Theorem}
\newtheorem{lemma}[theorem]{Lemma}
\newtheorem{proposition}[theorem]{Proposition}
\newtheorem{corollary}[theorem]{Corollary}
\theoremstyle{definition}
\newtheorem{definition}[theorem]{Definition}
\newtheorem{observation}[theorem]{Observation}
\newtheorem{remark}[theorem]{Remark}
\newtheorem{conjecture}[theorem]{Conjecture}

\newcommand{\Pq}{\Pi_q}
\newcommand{\PH}{P_H}
\newcommand{\eH}{e_H}
\newcommand{\Hq}{H_q(S_n)}
\newcommand{\Z}{\mathbb{Z}}
\newcommand{\Q}{\mathbb{Q}}
\newcommand{\tr}{\operatorname{tr}}
\newcommand{\bl}[1]{\overline{#1}}

\title{Multi-Eigenvalue Rosetta Stone for the Staircase Bott--Samelson
Element: Palindromy via Bar-Involution\\[4pt]
\large{Verification at $n\le 6$; integrality and palindromy proved;
positivity open}}
\author{Clio}
\date{2026-04-23}

\begin{document}
\maketitle

\begin{abstract}
Let $\Pq=\prod_{j=1}^{\binom{n}{2}}(1+T_{i_j})$ be the staircase
Bott--Samelson element of the Iwahori--Hecke algebra $\Hq$, with
$\mathbf{i}=(s_1)(s_2s_1)(s_3s_2s_1)\cdots$ a reduced word for
$w_0\in S_n$. For each irreducible $V_\lambda$ with $r:=\dim V_\lambda^H\ge 1$
(where $H=\langle s_1,s_3,\ldots\rangle$), the operator
$\Pq|_{V_\lambda}$ has rank $r$; let $e_\lambda^{(1)},\ldots,e_\lambda^{(r)}
\in\overline{\Q(q)}$ be its nonzero eigenvalues.

The \emph{Multi-Eigenvalue Rosetta Stone} conjecture asserts that for
every $k\in\{1,\ldots,r\}$, the elementary symmetric function
$\sigma_k(e_\lambda^{(1)},\ldots,e_\lambda^{(r)})$ factors as
$q^{a_{\lambda,k}}(1+q)^{b_{\lambda,k}}P_{\lambda,k}(q)$ with
$P_{\lambda,k}\in\Z_{\ge 0}[q]$ palindromic and $P_{\lambda,k}(0)=1$.

We prove three of the four ingredients and verify the full conjecture
computationally for $n\le 6$.
\begin{itemize}
\item \textbf{Integrality} (Theorem~\ref{thm:integrality}): each
  $\sigma_k\in\Z[q]$.
\item \textbf{Palindromy} (Theorem~\ref{thm:palindromy}):
  $\sigma_k(q^{-1})=q^{-kN}\sigma_k(q)$ with $N=\binom{n}{2}$,
  via the bar-involution $\bl{b_s}=q^{-1}b_s$ transported through
  exterior powers.
\item \textbf{Positivity on $q>0$} (Theorem~\ref{thm:realposivity}):
  each $\sigma_k(q)>0$ for real $q>0$, via the eigenvalue-positivity
  theorem.
\item \textbf{Non-negativity of integer coefficients}: open. Verified
  at $n\le 6$ for all $(\lambda,k)$.
\end{itemize}

We also exhibit a new datum at $n=6$: the trace of $\Pq$ on
$V_{(3,2,1)}$ is divisible by $(1+q)^1$ only --- not by
$(1+q)^{\lfloor n/2\rfloor}=(1+q)^3$. This refutes any direct
generalization of the trace-level $(1+q)^{\lfloor n/2\rfloor}$ result
(proved for $n\le 5$ in~\cite{clio2026milddlefactor}) to $n\ge 6$.
\end{abstract}

\section{Setup}

Let $\Hq$ denote the Iwahori--Hecke algebra of $S_n$ over $\Z[q,q^{-1}]$
with generators $T_1,\ldots,T_{n-1}$ and relations
\[
  T_i^2 = (q-1)T_i+q,\qquad
  T_iT_{i+1}T_i = T_{i+1}T_iT_{i+1},\qquad
  T_iT_j = T_jT_i\ (|i-j|\ge 2).
\]
Write $b_s:=1+T_s$. Then $b_s^2=(1+q)b_s$ (Lemma~\ref{lem:bs-idem}
below). Fix the staircase reduced expression
\[
  \mathbf{i} = (s_1)(s_2s_1)(s_3s_2s_1)\cdots(s_{n-1}\cdots s_1),
  \qquad \ell(\mathbf{i})=N:=\tbinom{n}{2},
\]
and define the staircase Bott--Samelson element
\[
  \Pq := \prod_{j=1}^{N} b_{i_j}
       = \prod_{j=1}^{N} (1+T_{i_j})\ \in\ \Hq.
\]
Let $H:=\langle s_1,s_3,s_5,\ldots\rangle\cong(S_2)^{\lfloor n/2\rfloor}$
be the parabolic generated by the odd Coxeter generators, and write
$\PH=\prod_{i\in\{1,3,5,\ldots\}\cap[1,n-1]}b_i$. The $H$-symmetrizer
$\eH:=\PH/(1+q)^{\lfloor n/2\rfloor}$ is an idempotent on any
$\Hq$-module, and $\operatorname{im}(\eH|_V)=V^H$.

The operator $\Pq|_{V_\lambda}$ has rank $r_\lambda:=\dim V_\lambda^H$
(proved for every $\lambda$ in~\cite{clio2026hinvariant}). Let
$e_\lambda^{(1)},\ldots,e_\lambda^{(r_\lambda)}\in\overline{\Q(q)}$
denote its nonzero eigenvalues and
\[
  \sigma_k := \sigma_k(e_\lambda^{(1)},\ldots,e_\lambda^{(r_\lambda)})
\]
the $k$th elementary symmetric function, $1\le k\le r_\lambda$.

\begin{lemma}\label{lem:bs-idem}
$b_s^2=(1+q)b_s$; hence $\PH^2=(1+q)^{\lfloor n/2\rfloor}\PH$ and $\eH$
is idempotent.
\end{lemma}
\begin{proof}
$(1+T_s)^2=1+2T_s+T_s^2=1+2T_s+(q-1)T_s+q=(1+q)(1+T_s)$.
Odd generators pairwise commute, so $\PH^2=\prod_i(1+T_i)^2=(1+q)^s\PH$
where $s=\lfloor n/2\rfloor$.
\end{proof}

\section{The Multi-Eigenvalue Rosetta Stone Conjecture}

\begin{conjecture}[Multi-Eigenvalue Rosetta Stone]\label{conj:main}
For every $n\ge 1$ and every irreducible $V_\lambda$ of $S_n$ with
$r_\lambda:=\dim V_\lambda^H\ge 1$, and for every
$k\in\{1,\ldots,r_\lambda\}$, there exist integers
$a_{\lambda,k},b_{\lambda,k}\ge 0$ and a polynomial
$P_{\lambda,k}\in\Z_{\ge 0}[q]$ with $P_{\lambda,k}(0)=1$ and
$P_{\lambda,k}$ palindromic, such that
\begin{equation}\label{eq:rosetta}
  \sigma_k = q^{a_{\lambda,k}}(1+q)^{b_{\lambda,k}}P_{\lambda,k}(q).
\end{equation}
\end{conjecture}

The conjecture decomposes into four structural claims:
\begin{enumerate}
\item[(I)]  \textbf{Integrality}: $\sigma_k\in\Z[q]$.
\item[(II)] \textbf{Palindromy}: after factoring out the largest $q^a$,
  the remainder $\sigma_k/q^a$ is a palindromic polynomial.
\item[(III)] \textbf{Non-negativity}: all coefficients of $\sigma_k$
  are non-negative.
\item[(IV)] \textbf{Normalization}: the lowest-order (and equivalently
  highest-order) nonzero coefficient of $\sigma_k$ is $1$.
\end{enumerate}
Given~(I)--(IV), \eqref{eq:rosetta} follows: set $a=\min\{j:\sigma_k[q^j]
\ne 0\}$, factor out the largest $(1+q)^b$ dividing $\sigma_k/q^a$
(palindromy divided by the palindrome $(1+q)^b$ remains
palindromic~\cite[standard]{stanley}), and observe $P(0)=1$ by~(IV).

\section{Integrality}

\begin{theorem}[Integrality]\label{thm:integrality}
For every $\lambda\vdash n$ and every $k\ge 0$,
$\sigma_k(e_\lambda^{(1)},\ldots,e_\lambda^{(r_\lambda)})\in\Z[q]$.
\end{theorem}

\begin{proof}
The irreducible $\Hq$-module $V_\lambda$ admits an integral form:
Dipper--James show that the cell module (Specht module) for the
Kazhdan--Lusztig left cell labeled by $\lambda$ is a free
$\Z[q,q^{-1}]$-module on which $\Hq$ acts by matrices in
$\operatorname{Mat}_{d\times d}(\Z[q,q^{-1}])$ (with $d=\dim
V_\lambda$). In fact, taking the integral Kazhdan--Lusztig basis,
$T_i$ and hence $b_i=1+T_i$ act by matrices with entries in
$\Z_{\ge 0}[q]$ (W-graph positivity, provable for type $A$ classically
and in general via Elias--Williamson).

The characteristic polynomial of any such matrix
$M\in\operatorname{Mat}_{d}(\Z[q,q^{-1}])$ lies in $\Z[q,q^{-1}][x]$.
For $M=\Pq|_{V_\lambda}$ the matrix entries are in $\Z_{\ge 0}[q]$
(product of matrices with entries in $\Z_{\ge 0}[q]$), so
$\det(xI-M)\in\Z[q][x]$.

The nontrivial factor of $\det(xI-M)$ is
$\tilde\chi_\lambda(x,q)=\prod_{i=1}^{r_\lambda}(x-e_\lambda^{(i)})$,
and its coefficients are (signed) elementary symmetric functions of the
$e_\lambda^{(i)}$. Each coefficient of $\det(xI-M)$ lies in $\Z[q]$;
hence $\sigma_k\in\Z[q]$.
\end{proof}

\section{Palindromy via Bar-Involution and Exterior Powers}

The key structural fact is that $\Pq$ is self-bar-dual up to a power of
$q^{-1}$; this transports through exterior powers to give palindromy of
each $\sigma_k$.

\begin{lemma}[Bar-transform of $b_s$]\label{lem:bs-bar}
The bar involution $\bl{\cdot}:\Hq\to\Hq$ defined by
$\bl{q}=q^{-1}$, $\bl{T_w}=T_{w^{-1}}^{-1}$ satisfies
$\bl{b_s}=q^{-1}b_s$.
\end{lemma}
\begin{proof}
From $T_s^2=(q-1)T_s+q$ we get $T_s^{-1}=q^{-1}(T_s-(q-1))=q^{-1}T_s+
q^{-1}-1$. Therefore
\[
  \bl{b_s}=1+\bl{T_s}=1+T_s^{-1}=q^{-1}T_s+q^{-1}=q^{-1}(T_s+1)=q^{-1}b_s.
\qedhere
\]
\end{proof}

\begin{lemma}[Bar-transform of $\Pq$]\label{lem:piq-bar}
$\bl{\Pq}=q^{-N}\Pq$, with $N=\binom{n}{2}$.
\end{lemma}
\begin{proof}
The bar involution is a $\Z$-algebra automorphism of $\Hq$. Hence
$\bl{\Pq}=\bl{\prod_j b_{i_j}}=\prod_j\bl{b_{i_j}}=\prod_j (q^{-1}
b_{i_j})=q^{-N}\prod_j b_{i_j}=q^{-N}\Pq$.
\end{proof}

\begin{theorem}[Palindromy]\label{thm:palindromy}
For every $\lambda\vdash n$ with $r_\lambda\ge 1$ and every
$k\in\{1,\ldots,r_\lambda\}$,
\begin{equation}\label{eq:palindromy}
  \sigma_k(q^{-1}) = q^{-kN}\,\sigma_k(q).
\end{equation}
Equivalently, writing $\sigma_k$ as a Laurent polynomial in $q$, the
coefficient of $q^j$ equals the coefficient of $q^{kN-j}$ for every
$j\in\Z$.
\end{theorem}

\begin{proof}
Let $M_\lambda$ be the matrix of $\Pq$ acting on $V_\lambda$ in the
integral Kazhdan--Lusztig basis (Dipper--James cell basis). This basis
is bar-invariant, so the bar involution of an entry
$(M_\lambda)_{ij}\in\Z[q]$ is $(M_\lambda)_{ij}(q^{-1})$. Combined with
Lemma~\ref{lem:piq-bar},
\begin{equation}\label{eq:M-bar}
  M_\lambda(q^{-1}) = q^{-N}\,M_\lambda(q).
\end{equation}

Fix $k\in\{1,\ldots,r_\lambda\}$. The $k$th exterior power of a linear
operator is a functor on matrices; it scales as
$\Lambda^k(cM)=c^k\Lambda^k(M)$ for any scalar $c$ (since
$\Lambda^k(cM)$ acts on $v_1\wedge\cdots\wedge v_k\mapsto
cMv_1\wedge\cdots\wedge cMv_k=c^k\Lambda^k(M)(v_1\wedge\cdots\wedge v_k)$).
Applying $\Lambda^k$ to~\eqref{eq:M-bar}:
\[
  \Lambda^k M_\lambda(q^{-1}) = q^{-kN}\,\Lambda^k M_\lambda(q).
\]
Taking traces and using
$\sigma_k=\tr(\Lambda^k M_\lambda)=\tr(\Lambda^k\Pq|_{V_\lambda})$
(basis-independent):
\[
  \sigma_k(q^{-1}) = q^{-kN}\,\sigma_k(q).\qedhere
\]
\end{proof}

\begin{corollary}[Palindromic structure]\label{cor:palstructure}
Write $\sigma_k=\sum_{j\ge 0}c_{k,j}q^j$. Then $c_{k,j}=c_{k,kN-j}$ for
every $j\ge 0$ (where $c_{k,j}=0$ if $j>kN$). In particular, if $a$ is
the smallest $j$ with $c_{k,j}\ne 0$ and $D$ is the largest, then
$a+D=kN$, and $\sigma_k/q^a$ is a palindromic polynomial of degree
$D-a=kN-2a$.
\end{corollary}
\begin{proof}
Directly from the identity $c_{k,j}=c_{k,kN-j}$ in
Theorem~\ref{thm:palindromy}.
\end{proof}

\begin{corollary}[Divisibility by $(1+q)^b$ preserves palindromy]\label{cor:palpres}
If $\sigma_k\in\Z[q]$ is palindromic in the sense of
Corollary~\ref{cor:palstructure} and $(1+q)^b$ divides $\sigma_k/q^a$
for the maximal such $a$, then the quotient
$\sigma_k/(q^a(1+q)^b)$ is itself palindromic.
\end{corollary}
\begin{proof}
$(1+q)^b$ is palindromic of degree $b$. If $R=(1+q)^bQ$ is palindromic
of degree $D'=D-a$, then $q^{D'}R(q^{-1})=R(q)$, i.e.,
$q^{D'}(1+q^{-1})^bQ(q^{-1})=q^{D'-b}(1+q)^bQ(q^{-1})=(1+q)^bQ(q)$;
cancelling $(1+q)^b$ gives $Q(q)=q^{D'-b}Q(q^{-1})$, so $Q$ is
palindromic of degree $D'-b$.
\end{proof}

\section{Positivity on real \boldmath$q>0$}

\begin{theorem}[Positivity on $q>0$]\label{thm:realposivity}
For every real $q>0$, $\sigma_k(q)>0$.
\end{theorem}
\begin{proof}
The Eigenvalue Positivity Theorem~\cite{clio2026eigenvaluepositivity}
gives: every nonzero eigenvalue of $\Pq|_{V_\lambda}$, as an element of
$\overline{\Q(q)}$, is positive for every specialization $q>0$. Since
each $\sigma_k$ is the $k$th elementary symmetric function of positive
real numbers (after specialization), $\sigma_k(q)>0$.
\end{proof}

Theorem~\ref{thm:realposivity} rules out positive real roots of
$\sigma_k$ (viewed as a polynomial in $q$). Combined with palindromy,
this rules out any root pair $\{q_0, q_0^{-1}\}$ of positive real roots.
However, palindromy plus positivity on $q>0$ is \emph{not} enough to
conclude non-negativity of integer coefficients: the palindromic
polynomial $q^4-q^3+q^2-q+1=\Phi_{10}(q)$ is positive on $q>0$ but has
negative coefficients.

\section{Computational verification for \boldmath$n\le 6$}

We computed $\Pq|_{V_\lambda}$ in Young's seminormal (Hoefsmit) form
for every irreducible $V_\lambda$ of $S_n$ for $n\in\{3,4,5,6\}$,
extracted the characteristic polynomial, factored out the zero
eigenvalues, and decomposed each $\sigma_k$ as
$q^a(1+q)^bP(q)$. Table~\ref{tab:data} records the factorization. In
every case $P\in\Z_{\ge 0}[q]$ is palindromic with $P(0)=1$.

Palindromy is verified algebraically (as a polynomial identity
$\sigma_k(q)=q^{kN}\sigma_k(q^{-1})$) and agrees with the bar-involution
proof of Theorem~\ref{thm:palindromy}.

\begin{longtable}{lllll}
\caption{Factorization $\sigma_k=q^a(1+q)^bP(q)$ for every irrep of
$S_n$, $n\le 6$, with $r_\lambda\ge 1$. $P$ is palindromic,
non-negative, $P(0)=1$.}\label{tab:data}\\
\toprule
$n$ & $\lambda$ & $k$ & $(a,b)$ & $P(q)$ \\
\midrule
\endfirsthead
\multicolumn{5}{c}{\tablename\ \thetable{} -- continued}\\
\toprule
$n$ & $\lambda$ & $k$ & $(a,b)$ & $P(q)$ \\
\midrule
\endhead
$3$ & $(3)$    & $1$ & $(0,3)$  & $1$ \\
$3$ & $(2,1)$  & $1$ & $(1,1)$  & $1$ \\
\midrule
$4$ & $(4)$    & $1$ & $(0,6)$  & $1$ \\
$4$ & $(3,1)$  & $1$ & $(1,2)$  & $q^2+4q+1$ \\
$4$ & $(2,2)$  & $1$ & $(2,2)$  & $1$ \\
\midrule
$5$ & $(5)$    & $1$ & $(0,10)$  & $1$ \\
$5$ & $(4,1)$  & $1$ & $(1,4)$   & $q^4+7q^3+17q^2+7q+1$ \\
$5$ & $(4,1)$  & $2$ & $(3,12)$  & $q^2+4q+1$ \\
$5$ & $(3,2)$  & $1$ & $(2,2)$   & $q^4+9q^3+19q^2+9q+1$ \\
$5$ & $(3,2)$  & $2$ & $(5,8)$   & $q^2+4q+1$ \\
$5$ & $(3,1,1)$ & $1$ & $(3,2)$  & $q^2+4q+1$ \\
$5$ & $(2,2,1)$ & $1$ & $(4,2)$  & $1$ \\
\midrule
$6$ & $(6)$    & $1$ & $(0,15)$   & $1$ \\
$6$ & $(5,1)$  & $1$ & $(1,7)$    & $q^6+9q^5+36q^4+70q^3+36q^2+9q+1$ \\
$6$ & $(5,1)$  & $2$ & $(3,18)$   & $q^6+12q^5+52q^4+88q^3+52q^2+12q+1$ \\
$6$ & $(4,2)$  & $1$ & $(2,3)$    & $q^8+13q^7+79q^6+245q^5+365q^4+245q^3+79q^2+13q+1$ \\
$6$ & $(4,2)$  & $2$ & $(5,10)$   & $(q^2+4q+1)(q^8+17q^7+103q^6+313q^5+473q^4+\ldots)$ \\
$6$ & $(4,2)$  & $3$ & $(9,21)$   & $(q^2+4q+1)^3$ \\
$6$ & $(4,1,1)$ & $1$ & $(3,3)$   & $q^6+12q^5+52q^4+88q^3+52q^2+12q+1$ \\
$6$ & $(3,3)$   & $1$ & $(3,3)$   & $q^6+12q^5+52q^4+88q^3+52q^2+12q+1$ \\
$6$ & $(3,2,1)$ & $1$ & $(4,1)$   & $(q^2+4q+1)(q^4+11q^3+21q^2+11q+1)$ \\
$6$ & $(3,2,1)$ & $2$ & $(9,2)$   & $q^{10}+16q^9+\cdots+16q+1$ \\
$6$ & $(2,2,2)$ & $1$ & $(6,3)$   & $1$ \\
\bottomrule
\end{longtable}

\begin{remark}[A new datum at $n=6$]\label{rem:new-datum}
At $n=6$, $\lambda=(3,2,1)$, $k=1$: the trace
$\tr_{V_{(3,2,1)}}(\Pq)=\sigma_1$ is divisible by $(1+q)^1$ only, not by
$(1+q)^3=(1+q)^{\lfloor n/2\rfloor}$. This refutes the naive
generalization of the result~\cite{clio2026milddlefactor} (which
established $(1+q)^{\lfloor n/2\rfloor}\mid\tr_{V_\lambda}(\Pq)$ for
$n\le 5$) to arbitrary $n$. In particular, the per-eigenvalue
$(1+q)^{\lfloor n/2\rfloor}$-divisibility conjecture fails at
$n=6,\lambda=(3,2,1)$, since $\sigma_1=\sum e_i$ divisible by
$(1+q)^{\lfloor n/2\rfloor}$ would follow from per-eigenvalue
divisibility. Conjecture~\ref{conj:main} remains consistent: it only
requires $b_{\lambda,k}\ge 0$, which is satisfied.
\end{remark}

\begin{remark}[The universal palindrome $q^2+4q+1$]
Two separate phenomena feature the palindrome $q^2+4q+1$:
\begin{itemize}
\item $P_{(3,1),n=4,k=1}=q^2+4q+1$ (the original ``Rosetta Stone''
  palindrome at $n=4$).
\item It appears as a factor of $\sigma_r$ at every rank-$r$ irrep we
  have checked. At $n=6,\lambda=(4,2),r=3$: $P_{(4,2),k=3}=(q^2+4q+1)^3$.
\end{itemize}
A branching-rule explanation would identify $(q^2+4q+1)$ as the
``atomic'' palindrome associated to $V_{(3,1)}$. We did not succeed in
formalizing such a branching rule: the obvious restriction $S_n\to
S_{n-1}$ is incompatible with the subgroup $H=\langle s_1,s_3,\ldots
\rangle$, since the top odd generator lies in $S_n\setminus S_{n-1}$.
\end{remark}

\section{Gaps and open problems}

Conjecture~\ref{conj:main} reduces (given integrality, palindromy, and
positivity on $q>0$) to:
\begin{enumerate}
\item[(G1)] \textbf{Non-negativity of integer coefficients.}
  Prove $\sigma_k\in\Z_{\ge 0}[q]$ for every $(\lambda,k)$.
\item[(G2)] \textbf{Normalization.} Prove the smallest nonzero
  coefficient of $\sigma_k$ is $1$.
\end{enumerate}

\subsection*{Attempted routes for (G1) and why they stalled}

\textbf{Route 1: Direct KL positivity.} The matrix of $\Pq$ in the KL
basis has entries in $\Z_{\ge 0}[q]$ (W-graph positivity). The trace
$\sigma_1$ is thus in $\Z_{\ge 0}[q]$, settling (G1) for $k=1$.
However for $k\ge 2$, $\sigma_k=\sum_{|S|=k}\det(M[S,S])$ is a sum of
$k\times k$ principal minors of $M$, and these determinants can have
mixed signs even when $M$ has non-negative entries. The direct sum is
not manifestly positive.

\textbf{Route 2: Lindström--Gessel--Viennot.} Using
$\Pq=b_{i_1}\cdots b_{i_N}$ in the KL basis, a telescoping application
of Cauchy--Binet gives
\[
  \sigma_k \;=\; \sum_{S_0=S_N,\,S_1,\ldots,S_{N-1}}\;
  \prod_{j=1}^{N}\det\bigl(b_{i_j}|_{S_{j-1},S_j}\bigr),
\]
where the sum runs over all closed sequences of $k$-subsets. If each
$b_{i_j}|_{V_\lambda}$ were a \emph{totally non-negative} matrix
(TNN), i.e., if every minor $\det(b_{i_j}|_{S,T})$ were non-negative,
this would give $\sigma_k\ge 0$ term-by-term.

We tested TNN-ness for $b_{s_1}|_{V_{(2,1)}}$ and
$b_{s_1}|_{V_{(3,1)}}$ in several natural bases (Kazhdan--Lusztig,
Murphy standard, Hoefsmit seminormal) and found: the diagonal
seminormal $b_{s_1}$ is TNN trivially (eigenvalue matrix); in the KL
basis the entries are non-negative but the 2-minors are not in general.
So TNN-ness does \emph{not} hold in the natural bases, and the LGV
route stalls.

\textbf{Route 3: Verdier duality / Soergel decomposition.} In the
Elias--Williamson Hecke category, $\Pq$ corresponds to the
Bott--Samelson object $BS(\mathbf{i})$, which decomposes as a direct
sum of indecomposable Soergel bimodules $BS(\mathbf{i})\simeq
\bigoplus_w m_w(\mathbf{i},q)B_w$ with Kazhdan--Lusztig multiplicities
$m_w\in\Z_{\ge 0}[q,q^{-1}]$. The character of each $B_w$ on $V_\lambda$
is a polynomial in $q$. This would yield an explicit Schur-positive-type
expansion of $\sigma_1=\tr_{V_\lambda}(\Pq)$; for $k\ge 2$,
$\sigma_k=\tr(\Lambda^k\Pq)$ does not come from an element of $\Hq$
(exterior power is not functorial on $\Hq$-modules), so the direct
Soergel route applies to $\sigma_1$ only.

\subsection*{What would close (G2)}

Normalization $P(0)=1$ is equivalent to: the lowest nonzero coefficient
of $\sigma_k$ equals $1$. From palindromy this coincides with the
leading coefficient. The leading coefficient of $\sigma_k$ is the
$q^{kN}$-coefficient of $\sigma_k$, which by palindromy equals the
constant term of $\sigma_k/q^a$ for $a=kN-\deg(\sigma_k)$. Computation
shows this constant term is $1$ in all observed cases; a structural
proof eludes us.

\section*{Acknowledgements}

Robin Langer for the conjecture formulation. Computational work used
Hoefsmit seminormal matrices over $\Q(q)$ with SymPy; the reference
implementation is
\texttt{/home/clio/projects/bms-test/idempotent\_eigenvalues.py} and
\texttt{n6\_multi\_eigenvalue.py}. Palindromy of the $n=6$ data was
verified in \texttt{verify\_palindromy.py}.

\begin{thebibliography}{9}
\bibitem{clio2026hinvariant} Clio,
\emph{The H-invariant theorem: $\operatorname{rank}(\Pq|_{V_\lambda})
=\dim V_\lambda^H$}. 2026-04-17.
\bibitem{clio2026eigenvaluepositivity} Clio,
\emph{Eigenvalue positivity for the staircase Bott--Samelson element}.
2026-04-17.
\bibitem{clio2026milddlefactor} Clio,
\emph{Middle factor of the Eigenvalue Rosetta Stone:
$(1+q)^{\lfloor n/2\rfloor}$-divisibility of $\tr_{V_\lambda}(\Pq)$}.
2026-04-23.
\bibitem{stanley} R. Stanley, \emph{Enumerative Combinatorics, Vol. 1},
Cambridge University Press, 2012.
\end{thebibliography}

\end{document}
